\subsection{Definition of measurable functions and sets}

DEFINITION 1. --- Let X be a locally compact space, \(\mu\) a measure on X. A mapping f of X into a topological space F is said to be measurable for the measure \(\mu\) (or to be \(\mu\) -measurable) if, for every compact subset K of X, there exist a \(\mu\) -negligible set \(N \subset K\) and a partition of K - N formed by a sequence (finite or infinite) ( \(K_{n}\) ) of compact sets, such that the restriction of f to each \(K_{n}\) is continuous.

It is clear that every continuous mapping of X into F is measurable.

Note that if \(\mu\) and \(\nu\) are two measures on X such that every \(\mu\) -negligible set is \(\nu\) -negligible, then every \(\mu\) -measurable function is also \(\nu\) -measurable (cf.~Ch. V, §5, Nos. 5, 6).

Definition 1 may be transformed into the following criterion:

PROPOSITION 1. --- For a mapping f of X into F to be measurable, it is necessary and sufficient that, for every compact set \(K \subset X\) and every number \(\varepsilon > 0\) , there exist a compact set \(K_{1} \subset K\) such that \(|\mu|(\mathrm{K} - \mathrm{K}_{1}) \leqslant \varepsilon\) and the restriction of f to \(K_{1}\) is continuous.

If this condition is fulfilled, we may define recursively a sequence of pairwise disjoint compact sets \(K_{n} \subset K\) such that \(|\mu|(\mathrm{K} - \bigcup_{i=1}^{n} \mathrm{K}_{i}) \leqslant 1/n\) and such that the restriction of f to each \(K_{n}\) is continuous (§4, No.~6, Th. 4); the complement with respect to K of the union of the \(K_{n}\) is then negligible (§4, No.~5, Cor. of Prop. 7), thus f is measurable. Conversely, suppose there exist a negligible set \(N \subset K\) and a partition \((\mathrm{K}_{n})\) of K - N formed of compact sets such that the restriction of f to each \(K_{n}\) is continuous; for every \(\varepsilon > 0\) there exists an integer n such that, if \(H = \bigcup_{i=1}^{n} K_{i}\) , then \(|\mu|(\mathrm{K} - \mathrm{H}) \leqslant \varepsilon\) (§4, No.~5, Cor. of Prop. 7); the set H is compact, the \(K_{i}\) ( \(1 \leqslant i \leqslant n\) ) form a finite partition of H into compact sets, and the restriction of f to each \(K_{i}\) is continuous; therefore the restriction of f to H is continuous.

PROPOSITION 2. --- Let \((\mathrm{F}_{n})\) be a sequence of topological spaces and, for each n, let \(f_{n}\) be a measurable mapping of X into \(F_{n}\) . For every compact set \(K \subset X\) and every \(\varepsilon > 0\) , there exists a compact set \(K_{0} \subset K\) such that \(|\mu|(\mathrm{K} - \mathrm{K}_{0}) \leqslant \varepsilon\) and the restriction to \(K_{0}\) of each of the functions \(f_{n}\) is continuous.

For each integer \(n \geqslant 1\) there exists a compact set \(K_{n} \subset K\) such that \(|\mu|(\mathrm{K} - \mathrm{K}_{n}) \leqslant \varepsilon/2^{n}\) and the restriction of \(f_{n}\) to \(K_{n}\) is continuous. The set \(K_{0} = \bigcap_{n=1}^{\infty} K_{n}\) is compact, the restrictions to \(K_{0}\) of all of the functions \(f_{n}\) are continuous and, since \(K - K_{0}\) is contained in the union of the \(K - K_{n}\) , \(|\mu|(\mathrm{K} - \mathrm{K}_{0}) \leqslant \sum_{n=1}^{\infty} \varepsilon/2^{n} = \varepsilon\) .

DEFINITION 2. --- A subset A of X is said to be measurable if its characteristic function \(\varphi_{A}\) is measurable.

In view of Def. 1, it comes to the same to say that a measurable set A is a set such that, for every compact set K, there exist a negligible set \(N \subset K\) and a partition \((K_{n})\) of K - N formed by a sequence of compact sets each of which is contained either in \(K \cap A\) or in \(K \cap C A\) .

This definition yields at once the following criterion:

PROPOSITION 3. --- For a set A to be measurable, it is necessary and sufficient that, for every compact set K, A∩K be integrable.

The condition is necessary, because the union of a sequence of integrable sets \(A_{n}\) is integrable when \(\sum_{n}|\mu|(\mathrm{A}_{n})\) is finite ( \(\S4\) , No.~5, Cor. of Prop. 8). The condition is sufficient because, for every integrable set B, there exist a negligible set \(N \subset B\) and a partition of B - N into a sequence of compact sets ( \(\S4\) , No.~6, Cor. 2 of Th. 4).

COROLLARY 1. --- The open sets and the closed sets are measurable.

In particular, the entire space \(\mathrm{X}\) is measurable.

COROLLARY 2. --- If X is metrizable, then every Souslin subset A of X (GT, IX, §6, No.~2) is μ-measurable for every measure μ on X.

By virtue of Prop. 3, it suffices to verify that every relatively compact Souslin set A is \(\mu\) -integrable. Now, such a set A is capacitable for \(|\mu|^*\) (GT, IX, §6, No.~9, Th. 5). Therefore, for every \(\varepsilon > 0\) there exists a compact subset K of A such that \(|\mu|^*(A) \leqslant |\mu|^*(K) + \varepsilon = |\mu|(K) + \varepsilon\) . Let U be a relatively compact open set in X containing A such that

\[
| \mu | (\mathrm{U}) = | \mu | ^ {*} (\mathrm{U}) \leqslant | \mu | ^ {*} (\mathrm{A}) + \varepsilon .
\]

Then \(|\mu|^*(\mathrm{U} - \mathrm{K}) = |\mu|(\mathrm{U}) - |\mu|(\mathrm{K}) \leqslant 2\varepsilon\) , therefore \(|\mu|^*(\mathrm{A} - \mathrm{K}) \leqslant 2\varepsilon\) , which proves that A is \(\mu\) -integrable (§4, No.~6, Cor. 1 of Th. 4).

